% it/figures/methodology/tab_workflow.tex — versione italiana
\begin{table}[!ht]
    \centering
    \small
    \renewcommand{\arraystretch}{1.3}
    \caption[Il workflow rilevare--tradurre--valutare.]{La pipeline della tesi. Ogni
    fase prende l'output della precedente come proprio input, così che l'analisi
    vincoli il design e il design sia ciò che la valutazione testa.}
    \label{tab:workflow}
    \begin{tabular}{@{}>{\raggedright\arraybackslash}p{2.3cm}>{\raggedright\arraybackslash}p{3.5cm}>{\raggedright\arraybackslash}p{3.2cm}>{\raggedright\arraybackslash}p{3.2cm}@{}}
        \toprule
        \textbf{Fase} & \textbf{Metodo} & \textbf{Input} & \textbf{Output} \\
        \midrule
        \textbf{Rilevare}
            & PCA, autoencoder, analisi delle anomalie e di preallarme
            & Osservazioni RAPID; rianalisi Copernicus
            & Struttura latente interpretabile del sistema \\
        \textbf{Tradurre}
            & Grammatica di design applicata alle caratteristiche recuperate; build basato su browser
            & Struttura latente da \emph{rilevare}
            & Tre condizioni rappresentazionali (statica, dinamica, esperienziale) \\
        \textbf{Valutare}
            & Studio a due gruppi; valutazione basata su questionario; analisi esplorativa
            & Rappresentazioni da \emph{tradurre}
            & Evidenza su come i non specialisti comprendono il sistema \\
        \bottomrule
    \end{tabular}
\end{table}
